% =================================================
% 文件: 物联网.tex
% 章节: 物联网技术
% 版本: v1.0
% =================================================

\documentclass[12pt,UTF8]{ctexbook}

\input{../common/page}

\xeCJKsetup{AutoFallBack=true}
\setCJKfamilyfont{hei}{SimHei}
\setCJKfamilyfont{kai}{KaiTi}

\usepackage{titletoc}
\titlecontents{chapter}[0pt]{\vspace{3mm}\bf\addvspace{2pt}\filright}{\contentspush{\thecontentslabel\hspace{0.8em}}}{}{\titlerule*[8pt]{.}\contentspage}

\usepackage[colorlinks,linkcolor=blue,citecolor=blue,urlcolor=blue]{hyperref}

\ctexset{
	part/name= {第,卷},
	part/number={\chinese{part}},
	chapter/name={第,章},
	chapter/number={\arabic{chapter}}
}

\usepackage{graphicx}
\graphicspath{{Images/}}

\usepackage[perpage]{footmisc}

\usepackage{enumitem}

\title{\heiti\zihao{0} 物联网技术入门指南}
\author{WangFei}
\date{\today}

\begin{document}

\maketitle
\tableofcontents

\frontmatter
\chapter{前言}

本指南面向初学者，详细介绍物联网的基础概念、架构、通信协议、传感器技术、平台开发等知识。

\mainmatter

\chapter{物联网基础概念}
\label{chap:iot_basics}

\section{什么是物联网}
\label{sec:what_is_iot}

物联网（IoT, Internet of Things）是将物理设备通过网络连接起来，实现数据采集、传输和处理的系统。

物联网的特点：
\begin{itemize}
    \item \textbf{互联性}：设备之间相互连接
    \item \textbf{智能化}：设备具备智能处理能力
    \item \textbf{感知能力}：设备可以感知环境
    \item \textbf{大规模}：大量设备同时接入
\end{itemize}

\section{物联网架构}
\label{sec:iot_architecture}

物联网三层架构：
\begin{table}[htbp]
    \centering
    \caption{物联网架构}
    \begin{tabular}{|c|p{8cm}|}
        \hline
        \textbf{层级} & \textbf{说明} \\
        \hline
        感知层 & 传感器、执行器、RFID \\
        \hline
        网络层 & 通信网络、网关、协议 \\
        \hline
        应用层 & 数据处理、平台、应用 \\
        \hline
    \end{tabular}
\end{table}

\section{物联网应用场景}
\label{sec:iot_applications}

物联网应用场景：
\begin{table}[htbp]
    \centering
    \caption{物联网应用场景}
    \begin{tabular}{|c|p{8cm}|}
        \hline
        \textbf{领域} & \textbf{应用} \\
        \hline
        智能家居 & 智能照明、安防监控、智能家电 \\
        \hline
        工业物联网 & 设备监控、生产自动化、预测维护 \\
        \hline
        智慧城市 & 智慧交通、环境监测、智能电网 \\
        \hline
        医疗健康 & 远程监护、可穿戴设备、智能医疗 \\
        \hline
        农业 & 智能灌溉、环境监测、精准农业 \\
        \hline
    \end{tabular}
\end{table}

\chapter{传感器技术}
\label{chap:sensors}

\section{传感器基础}
\label{sec:sensor_basics}

传感器是物联网的感知器官，将物理量转换为电信号。

传感器分类：
\begin{table}[htbp]
    \centering
    \caption{传感器分类}
    \begin{tabular}{|c|p{8cm}|}
        \hline
        \textbf{类型} & \textbf{说明} \\
        \hline
        温度传感器 & 测量温度 \\
        \hline
        湿度传感器 & 测量湿度 \\
        \hline
        压力传感器 & 测量压力 \\
        \hline
        光照传感器 & 测量光照强度 \\
        \hline
        加速度传感器 & 测量运动状态 \\
        \hline
        气体传感器 & 检测气体浓度 \\
        \hline
    \end{tabular}
\end{table}

\section{常见传感器}
\label{sec:common_sensors}

常见传感器介绍：
\begin{table}[htbp]
    \centering
    \caption{常见传感器}
    \begin{tabular}{|c|c|p{5cm}|}
        \hline
        \textbf{传感器} & \textbf{接口} & \textbf{用途} \\
        \hline
        DHT11 & 单总线 & 温湿度检测 \\
        \hline
        DS18B20 & 单总线 & 温度检测 \\
        \hline
        HC-SR04 & 超声波 & 距离测量 \\
        \hline
        MQ-2 & 模拟 & 可燃气体检测 \\
        \hline
        BH1750 & I2C & 光照强度 \\
        \hline
    \end{tabular}
\end{table}

\chapter{通信协议}
\label{chap:protocols}

\section{物联网通信协议}
\label{sec:iot_protocols}

物联网通信协议分为有线和无线：

无线协议：
\begin{table}[htbp]
    \centering
    \caption{无线通信协议}
    \begin{tabular}{|c|c|p{5cm}|}
        \hline
        \textbf{协议} & \textbf{范围} & \textbf{用途} \\
        \hline
        Wi-Fi & 中距离 & 智能家居、网关 \\
        \hline
        Bluetooth & 短距离 & 可穿戴设备 \\
        \hline
        Zigbee & 中距离 & 工业、智能家居 \\
        \hline
        LoRa & 远距离 & 物联网网关 \\
        \hline
        NB-IoT & 远距离 & 低功耗物联网 \\
        \hline
        5G & 高速 & 工业物联网 \\
        \hline
    \end{tabular}
\end{table}

有线协议：
\begin{table}[htbp]
    \centering
    \caption{有线通信协议}
    \begin{tabular}{|c|p{8cm}|}
        \hline
        \textbf{协议} & \textbf{说明} \\
        \hline
        RS-232 & 短距离串口通信 \\
        \hline
        RS-485 & 长距离串口通信 \\
        \hline
        CAN & 车载网络 \\
        \hline
        Ethernet & 局域网通信 \\
        \hline
    \end{tabular}
\end{table}

\section{MQTT}
\label{sec:mqtt}

MQTT（Message Queuing Telemetry Transport）是轻量级消息协议：

MQTT 特点：
\begin{itemize}
    \item 轻量级
    \item 低功耗
    \item 发布/订阅模式
    \item QoS 保证
\end{itemize}

MQTT 服务质量（QoS）：
\begin{table}[htbp]
    \centering
    \caption{MQTT QoS}
    \begin{tabular}{|c|p{8cm}|}
        \hline
        \textbf{级别} & \textbf{说明} \\
        \hline
        QoS 0 & 最多一次，不保证送达 \\
        \hline
        QoS 1 & 至少一次，保证送达 \\
        \hline
        QoS 2 & 精确一次，保证只送达一次 \\
        \hline
    \end{tabular}
\end{table}

安装 MQTT 服务器（Mosquitto）：
\begin{verbatim}
dnf install -y mosquitto
systemctl enable --now mosquitto
\end{verbatim}

发布消息：
\begin{verbatim}
mosquitto_pub -h localhost -t "sensor/temperature" -m "25.5"
\end{verbatim}

订阅消息：
\begin{verbatim}
mosquitto_sub -h localhost -t "sensor/#"
\end{verbatim}

\section{CoAP}
\label{sec:coap}

CoAP（Constrained Application Protocol）是 RESTful 协议：

CoAP 特点：
\begin{itemize}
    \item 基于 UDP
    \item 轻量级
    \item RESTful 接口
    \item 资源发现
\end{itemize}

\chapter{物联网平台}
\label{chap:platforms}

\section{开源物联网平台}
\label{sec:open_source_platforms}

开源物联网平台：
\begin{table}[htbp]
    \centering
    \caption{开源物联网平台}
    \begin{tabular}{|c|p{8cm}|}
        \hline
        \textbf{平台} & \textbf{说明} \\
        \hline
        Eclipse IoT & 开源物联网框架 \\
        \hline
        ThingsBoard & 物联网数据平台 \\
        \hline
        OpenHAB & 智能家居平台 \\
        \hline
        Home Assistant & 开源智能家居 \\
        \hline
        Node-RED & 物联网流程编排 \\
        \hline
    \end{tabular}
\end{table}

安装 Home Assistant：
\begin{verbatim}
docker run -d \
  --name homeassistant \
  --privileged \
  --restart=unless-stopped \
  -e TZ=Asia/Shanghai \
  -v /path/to/config:/config \
  --network=host \
  ghcr.io/home-assistant/home-assistant:stable
\end{verbatim}

\section{云物联网平台}
\label{sec:cloud_platforms}

云厂商物联网平台：
\begin{table}[htbp]
    \centering
    \caption{云物联网平台}
    \begin{tabular}{|c|p{8cm}|}
        \hline
        \textbf{平台} & \textbf{说明} \\
        \hline
        AWS IoT & 亚马逊物联网服务 \\
        \hline
        Azure IoT Hub & 微软物联网服务 \\
        \hline
       阿里云 IoT & 阿里巴巴物联网平台 \\
        \hline
       腾讯云 IoT & 腾讯物联网平台 \\
        \hline
    \end{tabular}
\end{table}

\chapter{物联网设备开发}
\label{chap:device_development}

\section{开发板}
\label{sec:development_boards}

常用物联网开发板：
\begin{table}[htbp]
    \centering
    \caption{开发板}
    \begin{tabular}{|c|p{8cm}|}
        \hline
        \textbf{开发板} & \textbf{说明} \\
        \hline
        Arduino & 入门级开发板 \\
        \hline
        ESP32 & Wi-Fi+蓝牙，性价比高 \\
        \hline
        Raspberry Pi & 微型电脑 \\
        \hline
        STM32 & 工业级 MCU \\
        \hline
    \end{tabular}
\end{table}

\section{ESP32 开发}
\label{sec:esp32_development}

ESP32 是广泛使用的物联网芯片：

安装 ESP-IDF：
\begin{verbatim}
git clone --recursive https://github.com/espressif/esp-idf.git
cd esp-idf
./install.sh
. ./export.sh
\end{verbatim}

创建项目：
\begin{verbatim}
idf.py create-project myproject
cd myproject
idf.py set-target esp32
idf.py build
idf.py flash monitor
\end{verbatim}

ESP32 连接 Wi-Fi：
\begin{verbatim}
#include "esp_wifi.h"

void wifi_init() {
    wifi_config_t wifi_config = {
        .sta = {
            .ssid = "SSID",
            .password = "PASSWORD",
        },
    };
    esp_wifi_set_config(WIFI_IF_STA, &wifi_config);
    esp_wifi_start();
}
\end{verbatim}

\section{Raspberry Pi 开发}
\label{sec:raspberry_pi_development}

Raspberry Pi 是微型计算机：

安装操作系统：
\begin{verbatim}
# 使用 Raspberry Pi Imager
# 或手动烧录
dd if=raspios.img of=/dev/sdX bs=4M
\end{verbatim}

启用 SSH：
\begin{verbatim}
touch /boot/ssh
\end{verbatim}

Python 控制 GPIO：
\begin{verbatim}
import RPi.GPIO as GPIO
GPIO.setmode(GPIO.BCM)
GPIO.setup(18, GPIO.OUT)
GPIO.output(18, GPIO.HIGH)
\end{verbatim}

\chapter{物联网安全}
\label{chap:iot_security}

\section{安全挑战}
\label{sec:security_challenges}

物联网安全挑战：
\begin{itemize}
    \item \textbf{设备安全}：设备资源受限，安全能力弱
    \item \textbf{网络安全}：通信协议漏洞
    \item \textbf{数据安全}：数据泄露风险
    \item \textbf{身份认证}：设备认证困难
\end{itemize}

\section{安全措施}
\label{sec:security_measures}

物联网安全措施：
\begin{itemize}
    \item \textbf{加密通信}：使用 TLS/SSL
    \item \textbf{身份认证}：设备证书
    \item \textbf{固件更新}：安全升级机制
    \item \textbf{访问控制}：最小权限原则
\end{itemize}

\chapter{物联网数据处理}
\label{chap:data_processing}

\section{数据采集}
\label{sec:data_collection}

数据采集方法：
\begin{itemize}
    \item \textbf{定时采集}：定期采集传感器数据
    \item \textbf{事件触发}：根据事件采集
    \item \textbf{阈值触发}：超过阈值时采集
\end{itemize}

数据格式：
\begin{table}[htbp]
    \centering
    \caption{数据格式}
    \begin{tabular}{|c|p{8cm}|}
        \hline
        \textbf{格式} & \textbf{说明} \\
        \hline
        JSON & 轻量级数据交换格式 \\
        \hline
        CSV & 表格数据格式 \\
        \hline
        Protocol Buffers & 高效二进制格式 \\
        \hline
        MessagePack & 高效二进制格式 \\
        \hline
    \end{tabular}
\end{table}

\section{边缘计算}
\label{sec:edge_computing}

边缘计算是在设备端进行数据处理：

边缘计算优势：
\begin{itemize}
    \item 低延迟
    \item 减少网络流量
    \item 数据隐私保护
    \item 离线运行能力
\end{itemize}

边缘计算框架：
\begin{table}[htbp]
    \centering
    \caption{边缘计算框架}
    \begin{tabular}{|c|p{8cm}|}
        \hline
        \textbf{框架} & \textbf{说明} \\
        \hline
        OpenFog & 雾计算标准 \\
        \hline
        Apache Kafka Edge & 边缘数据处理 \\
        \hline
        TensorFlow Lite & 边缘机器学习 \\
        \hline
    \end{tabular}
\end{table}

\chapter{物联网实践}
\label{chap:practice}

\section{智能家居项目}
\label{sec:smart_home_project}

创建智能家居项目：
\begin{enumerate}
    \item 硬件选择：ESP32 + 传感器
    \item 通信协议：MQTT
    \item 平台：Home Assistant
    \item 功能：温湿度监测、智能开关
\end{enumerate}

项目代码结构：
\begin{verbatim}
smart-home/
├── sensors/          # 传感器驱动
├── mqtt/             # MQTT 客户端
├── config/           # 配置文件
└── main.py           # 主程序
\end{verbatim}

\section{环境监测项目}
\label{sec:environmental_monitoring}

环境监测项目：
\begin{enumerate}
    \item 传感器：温湿度、空气质量、光照
    \item 通信：LoRa 或 NB-IoT
    \item 数据存储：InfluxDB
    \item 可视化：Grafana
\end{enumerate}

\backmatter

\end{document}